\chapter{Installation und erste Rechnung}\label{sec:app-start}
\index{Anwendung}\index{Installation}
Dieses Benutzerhandbuch beschreibt Version \ToolboxVersion\ am angegebenen
Quellcommit. Es umfasst den integrierten M8-Funktionsstand und die vorhandenen
M9-Korrekturen. API-Bezeichner bleiben englisch. Der Paketimport heißt
\path{fem_toolbox}; die geplante Femantis-Umbenennung ist kein Bestandteil
dieses Quellstands. Die mathematischen Konventionen stehen im Theorieband;
Installation und wiederholbare Arbeitsschritte stehen hier.

\section{Support und optionale Werkzeuge}
\begin{table}[htbp]\centering
\begin{tabular}{p{.25\linewidth}p{.64\linewidth}}
\toprule
Umgebung & Umfang und Nachweis\\\midrule
CPython 3.11--3.13 & Unterstützter Paketbereich; Linux-CI prüft Kern,
Wheel und sdist für jede Version.\\
macOS arm64, 3.13 & Lokaler Paketnachweis mit frisch installiertem Wheel;
kein Versprechen für alle macOS-/Python-Kombinationen.\\
Windows & Keine eigenständige CI-Matrix. Die angegebenen Interpreterpfade
sind Installationshinweise, kein behaupteter Windows-Nachweis.\\
NumPy/SciPy/meshio & Kernabhängigkeiten; Mindestversionen und Python-Marker
stehen verbindlich in \path{pyproject.toml}.\\
Matplotlib & Extra \path{visualization}, Version $\geq3.9,<4$;
JPG-Rezepte laufen ohne grafische Sitzung.\\
PETSc & Optionaler separater nativer Linux-Runtime, real/complex;
SciPy bleibt Standard und Model/TOML verwenden SciPy.\\
Gmsh/ParaView & Externe Programme zur Netzvorbereitung/Ansicht.
Die mitgelieferten Netze brauchen keine Gmsh-Installation.\\
\bottomrule
\end{tabular}
\caption{Supportumfang des Handbuchkandidaten, keine pauschale Plattformfreigabe.}
\end{table}
Der Kern benötigt kein TeX. PDF-Erstellung ist ein getrenntes Docker-Werkzeug
mit fixierter Toolchain; \cite{latex-project} beschreibt das Dokumentensystem.
Die Extras \path{test} und \path{dev} sind Entwicklungswerkzeuge, keine
Voraussetzung für die folgenden Anwendungsrezepte. PETSc ist bewusst kein
pip-Extra: native Bibliothek, Scalar-Typ und ABI müssen zusammenpassen.

\section{Artefaktinstallation statt Quellbaumimport}
Wheel und Handbuch-Quellenarchiv müssen zum gleichen geprüften Commit gehören.
Der Paketstand wird über GitHub-Release-/CI-Artefakte bzw. lokal mit
\path{python -m build} gebaut; das Handbuch behauptet keine Veröffentlichung
von M9 auf PyPI. Nach dem Entpacken heißt das Verzeichnis \path{handbook}.
Die Beispielprogramme sind Teil dieses Quellenarchivs, nicht des Wheels.
In einem davon unabhängigen Arbeitsverzeichnis:
\begin{lstlisting}
python3.13 -m venv /tmp/fem-user-env
/tmp/fem-user-env/bin/python -m pip install /ABS/PATH/fem_toolbox-1.4.0-py3-none-any.whl
/tmp/fem-user-env/bin/python -m pip check
/tmp/fem-user-env/bin/python -I -c "import fem_toolbox; print(fem_toolbox.__file__)"
cd /ABS/PATH/handbook
/tmp/fem-user-env/bin/python -I examples/steady_heat.py
\end{lstlisting}
Die beiden absoluten Pfade ersetzen Sie durch Ihre Artefaktablage. Der
Import muss unter \path{site-packages} der neuen Umgebung liegen, nicht unter
\path{src}. Unter Windows lautet der Interpreterpfad beispielsweise
\path{C:/work/fem-user-env/Scripts/python.exe}; die Erstellung verwendet
\path{py -3.13 -m venv C:/work/fem-user-env}. Im Rezeptkatalog bedeutet
\path{python} immer genau diesen Interpreter. Optional kann die Umgebung
aktiviert werden. \path{-I} verhindert versehentlichen Import aus dem
Arbeitsverzeichnis oder \path{PYTHONPATH}; bestehende Benutzereinstellungen
ersetzen keinen isolierten Installationsnachweis.

Für JPGs installiert man in derselben Umgebung zusätzlich das zum Paketstand
gehörige Wheel mit Extra, nicht eine andere Paketversion:
\begin{lstlisting}
python -m pip install "/ABS/PATH/fem_toolbox-1.4.0-py3-none-any.whl[visualization]"
python -m pip check
python -I examples/field_visualization/run.py --output results/images
\end{lstlisting}
Auf einem Entwicklungscheckout wird dagegen nach \path{INSTALL.md} eine
editierbare Installation eingerichtet. Sie ist zweckmäßig für Änderungen,
aber nicht der Nachweis, dass ein ausgeliefertes Wheel seine eigenen Module
enthält. Für exakte Wiederholung zusätzlich Python-/Dependency-Versionen mit
\path{python -m pip freeze} erfassen. Hashes sichern Dateiintegrität; sie sind
keine digitale Signatur und keine Authentizitätsprüfung.

\section{Erster Erfolg und was er bedeutet}
Rezept \ref{recipe:steady-heat} benötigt keine Eingabedatei: Der Python-Code
baut sechs Tet4 im Quader $2\times3\times4$ m. Links gelten 400 K, rechts
Konvektion mit $h=4$ W/(m$^2$ K) gegen 300 K, Leitfähigkeit
$k=2$ W/(m K). Erwartet werden 320--400 K und 960 W Abfluss;
das Skript prüft die affine Lösung und die Leistungsbilanz vor dem Export.
\path{steady_heat.vtu} entsteht im aktuellen Arbeitsverzeichnis.

Öffnen Sie die Datei in ParaView, drücken Sie Apply, wählen Sie
\path{temperature_K} unter Point Data und setzen Sie den Farbbereich auf
320--400. \path{heat_flux_W_m2} ist rohe Volumenzellinformation; die
Triangle3-Zellzeilen enthalten Exportfüllwerte. Das affine Referenzfeld und
die Bilanz sind ein numerischer Verifikationsfall. Ein plausibles Farbbild
oder ein kleines Residuum belegen weder Netzkonvergenz eines neuen Modells
noch Übereinstimmung mit unabhängigen Experimenten.
